\chapter{Software Design}

\section{Introduction}
This chapter specifies the software architecture and detailed design of the Quantum Virtual Machine (QVM). Building an execution runtime for quantum algorithms presents an interesting dual challenge: the mathematics of quantum mechanics is strict and unforgiving, yet the software executing it must remain fast, modular, and easy to maintain on classical computers.

To address this, we designed QVM as a layered, object-oriented pipeline. The design separates high-level language parsing from physical hardware topology routing, and keeps quantum state mathematics isolated from classical control flow logic. In this chapter, we document our architectural choices, structural models, behavioral workflows, and data structures. We provide complete diagrams for every dimension of the system, including context boundaries, pipeline architecture, activity flows, class structures, sequence lifecycles, and data flows.

\section{Design Methodology}
The architecture of QVM was guided by five core software engineering principles:

\textbf{Separation of Concerns:} Each stage of the quantum program lifecycle is isolated inside its own module. The parser only cares about grammar rules and syntax trees; the decomposer only cares about gate identities; the transpiler only cares about coupling graphs and SWAP heuristics; and the simulator only cares about unitary linear algebra and measurement collapse. This isolation guarantees that extending one module (such as adding support for an ion-trap topology) never destabilizes another (such as the OpenQASM 3.0 parser).

\textbf{Pipeline Architecture Pattern:} The execution runtime follows a sequential data pipeline:
$$\text{Input Source} \longrightarrow \text{Parser} \longrightarrow \text{Circuit IR} \longrightarrow \text{Decomposer} \longrightarrow \text{Transpiler} \longrightarrow \text{Simulator} \longrightarrow \text{Output}$$
Information flows predictably in one direction. Each stage accepts a structured representation, applies a deterministic transformation, and emits a refined representation for the next stage.

\textbf{Strategy Pattern for Routing:} The transpiler implements the Strategy design pattern for qubit routing. Users can switch dynamically between Greedy Breadth-First Search (BFS) for rapid local compilation and the SABRE lookahead heuristic \cite{li2019sabre} for deep circuit depth reduction, without modifying client scripts or driver code.

\textbf{Hardware Abstraction through Intermediate Representation:} The central \texttt{QuantumCircuit} class establishes a hardware-neutral intermediate representation. By converting diverse input formats (OpenQASM 2.0, OpenQASM 3.0, JSON) into this common IR before performing decomposition or routing, QVM decouples user algorithms from target chip constraints, realizing true ``Write Once, Run Anywhere'' portability.

\textbf{Stateless Execution Engine:} Quantum simulations are stateless by default. Each circuit execution accepts an initial configuration, evolves the quantum state in memory, and returns the final amplitudes and measurement samples. This keeps the execution engine thread-safe, scalable across web endpoints, and trivial to test.

\section{System Overview}
Figure \ref{fig:context} illustrates the system context boundary of the Quantum Virtual Machine. It captures how external actors (quantum algorithm developers, students, researchers) and external tools (local file systems, cloud databases, and third-party quantum ecosystems like IBM Quantum) interact with the QVM core runtime.

\begin{figure}[!htb]
\centering
\resizebox{0.95\textwidth}{!}{%
\begin{tikzpicture}[
    actorbox/.style={rectangle, draw=blue!80!black, fill=blue!8, rounded corners=3pt, line width=1pt, align=center, font=\small\bfseries, minimum height=1cm, inner sep=6pt},
    systembox/.style={rectangle, draw=blue!90!black, fill=blue!2, rounded corners=6pt, line width=1.5pt, inner sep=12pt},
    subbox/.style={rectangle, draw=black!70, fill=white, rounded corners=3pt, line width=0.8pt, align=center, font=\footnotesize, inner sep=5pt},
    extbox/.style={rectangle, draw=orange!80!black, fill=orange!8, rounded corners=3pt, line width=1pt, align=center, font=\small\bfseries, minimum height=1cm, inner sep=6pt},
    flow/.style={-{Stealth[scale=1.0]}, line width=1.1pt, draw=black!75},
    biflow/.style={{Stealth[scale=1.0]}-{Stealth[scale=1.0]}, line width=1.1pt, draw=black!75}
]

% External Actors (Left)
\node[actorbox] (dev) at (-6.5, 2.5) {Quantum Algorithm\\Developer};
\node[actorbox] (student) at (-6.5, 0.0) {Student /\\Educator};
\node[actorbox] (researcher) at (-6.5, -2.5) {Quantum\\Researcher};

% System Boundary (Center)
\node[systembox, minimum width=6.5cm, minimum height=7.5cm] (qvm) at (0, 0) {};
\node[font=\bfseries\normalsize, text=blue!90!black, anchor=north] at (qvm.north) {Quantum Virtual Machine (QVM) Boundary};

\node[subbox, minimum width=5.5cm] (ifaces) at (0, 2.2) {\textbf{User Interfaces}\\CLI (\texttt{cli.py}) $\bullet$ FastAPI REST $\bullet$ Next.js GUI};
\node[subbox, minimum width=5.5cm] (compiler) at (0, 0.5) {\textbf{Compiler \& Transpiler Subsystem}\\Lark Parser $\bullet$ \texttt{QuantumCircuit} IR $\bullet$ SABRE Router};
\node[subbox, minimum width=5.5cm] (sim) at (0, -1.2) {\textbf{Dual Simulation Engines}\\Dense Statevector Engine $\bullet$ MPS Tensor Network};
\node[subbox, minimum width=5.5cm] (noise) at (0, -2.7) {\textbf{Open-System Noise Engine}\\Kraus Channels $\bullet$ Emulated Device Calibrations};

\draw[-{Stealth[scale=0.9]}, thick] (ifaces) -- (compiler);
\draw[-{Stealth[scale=0.9]}, thick] (compiler) -- (sim);
\draw[-{Stealth[scale=0.9]}, thick] (sim) -- (noise);

% External Systems (Right)
\node[extbox] (fs) at (6.5, 2.3) {Local File System\\(\texttt{.qasm}, \texttt{.json}, \texttt{.png})};
\node[extbox] (cloud) at (6.5, 0.0) {Supabase Cloud DB\\(Circuits \& History)};
\node[extbox] (ecosystem) at (6.5, -2.3) {External Ecosystem\\(IBM Quantum / Qiskit)};

% Connections
\draw[flow] (dev) -- node[above, font=\tiny\bfseries]{Run CLI Commands} (ifaces.west |- dev);
\draw[flow] (student) -- node[above, font=\tiny\bfseries]{Web GUI Interactions} (ifaces.west |- student);
\draw[flow] (researcher) -- node[above, font=\tiny\bfseries]{API / Benchmarks} (ifaces.west |- researcher);

\draw[biflow] (ifaces.east |- fs) -- node[above, font=\tiny\bfseries]{Read / Save Files} (fs);
\draw[biflow] (ifaces.east |- cloud) -- node[above, font=\tiny\bfseries]{Sync Circuits} (cloud);
\draw[flow] (compiler.east |- ecosystem) -- node[above, font=\tiny\bfseries]{Export OpenQASM} (ecosystem);

\end{tikzpicture}%
}
\caption{QVM System Context Diagram}
\label{fig:context}
\end{figure}

The context diagram establishes clear operational boundaries. Users supply circuit definitions through three interaction channels (terminal CLI, web dashboard, or REST API). The system validates the input, executes the transpilation and simulation pipeline, and exports results as numerical statevectors, measurement histograms, formatted JSON payloads, or standard OpenQASM strings compatible with external platforms like IBM Quantum.

\section{Architectural Design}
The internal organization of QVM follows a \textbf{7-Layer Pipeline Architecture}, illustrated in Figure \ref{fig:architecture}. This replaces bloated multi-page diagrams with a unified, publication-quality vector specification.

\begin{figure}[!htb]
\centering
\resizebox{0.92\textwidth}{!}{%
\begin{tikzpicture}[
  layer/.style={rectangle, draw=blue!80!black, fill=blue!5, rounded corners=4pt, minimum width=13.5cm, minimum height=1.15cm, font=\small, align=center, line width=0.9pt},
  flowarrow/.style={-{Stealth[scale=1.1]}, line width=1.2pt, draw=blue!85!black}
]

% Layer 1
\node[layer, fill=cyan!8, draw=cyan!80!black] (L1) at (0, 0) {
  \textbf{Layer 1: Input \& Interfaces}\\[1pt]
  \footnotesize CLI (\texttt{cli.py}) \quad$\bullet$\quad FastAPI REST Service (\texttt{api/app.py}) \quad$\bullet$\quad Next.js 16 Web Dashboard \quad$\bullet$\quad File Loader (.qasm, .json)
};

% Layer 2
\node[layer, fill=teal!8, draw=teal!80!black] (L2) at (0, -1.6) {
  \textbf{Layer 2: Parsing \& AST Validation}\\[1pt]
  \footnotesize Lark LALR(1) Parser (\texttt{qasm3.lark}) \quad$\bullet$\quad OpenQASM 2.0 Parser \quad$\bullet$\quad AST Generator \quad$\bullet$\quad Semantic Validator
};

% Layer 3
\node[layer, fill=blue!8, draw=blue!80!black] (L3) at (0, -3.2) {
  \textbf{Layer 3: Intermediate Representation (IR)}\\[1pt]
  \footnotesize \texttt{QuantumCircuit} Class (\texttt{ir.py}) \quad$\bullet$\quad Sequential Gate Stream \quad$\bullet$\quad Qubit Register Map \quad$\bullet$\quad Classical Bit Registers
};

% Layer 4
\node[layer, fill=indigo!8, draw=indigo!80!black] (L4) at (0, -4.8) {
  \textbf{Layer 4: Gate Decomposition}\\[1pt]
  \footnotesize Toffoli Synthesis ($CCX \to 6CX + 9U_1$) \cite{nielsen2010} \quad$\bullet$\quad Multi-Controlled Lowering (\texttt{synthesis.py}) \quad$\bullet$\quad Native Basis $\{U, CX\}$
};

% Layer 5
\node[layer, fill=orange!8, draw=orange!80!black] (L5) at (0, -6.4) {
  \textbf{Layer 5: Topology-Aware Transpiler \& Routing}\\[1pt]
  \footnotesize Coupling Topologies (Linear, Grid, Heavy-Hex) \quad$\bullet$\quad Greedy BFS Routing \quad$\bullet$\quad SABRE Lookahead \cite{li2019sabre} \quad$\bullet$\quad SWAP Inserter
};

% Layer 6
\node[layer, fill=purple!8, draw=purple!80!black] (L6) at (0, -8.0) {
  \textbf{Layer 6: Dual Simulation Engines \& Noise Model}\\[1pt]
  \footnotesize Statevector ($n \le 12$, NumPy tensor slicing) \quad$\bullet$\quad MPS ($n \ge 20$, SVD bond truncation $\chi$) \quad$\bullet$\quad Kraus Noise ($T_1, T_2$, Depol, Readout)
};

% Layer 7
\node[layer, fill=green!8, draw=green!80!black] (L7) at (0, -9.6) {
  \textbf{Layer 7: Output, Visualization \& Persistence}\\[1pt]
  \footnotesize State Amplitudes $|\psi\rangle$ \quad$\bullet$\quad Measurement Probabilities \quad$\bullet$\quad Matplotlib \& SVG Schematics \quad$\bullet$\quad Supabase Cloud Storage
};

% Connecting Arrows with Labels
\draw[flowarrow] (L1) -- node[right, font=\tiny\bfseries, text=black!70]{Raw Text / JSON} (L2);
\draw[flowarrow] (L2) -- node[right, font=\tiny\bfseries, text=black!70]{Validated AST} (L3);
\draw[flowarrow] (L3) -- node[right, font=\tiny\bfseries, text=black!70]{Circuit IR} (L4);
\draw[flowarrow] (L4) -- node[right, font=\tiny\bfseries, text=black!70]{Decomposed Circuit} (L5);
\draw[flowarrow] (L5) -- node[right, font=\tiny\bfseries, text=black!70]{Physical Circuit + SWAPs} (L6);
\draw[flowarrow] (L6) -- node[right, font=\tiny\bfseries, text=black!70]{Statevector / Counts} (L7);

\end{tikzpicture}%
}
\caption{QVM 7-Layer Modular Pipeline Architecture}
\label{fig:architecture}
\end{figure}

The responsibilities of each architectural layer are defined as follows:

\textbf{Layer 1 - Input Layer:} Ingests quantum circuits across three file formats (OpenQASM 3.0, OpenQASM 2.0, JSON) and three interface channels (CLI, REST API, Web GUI).

\textbf{Layer 2 - Parser Layer:} Uses Lark's LALR(1) parsing engine \cite{lark2023} driven by a formal EBNF grammar to build an Abstract Syntax Tree (AST). It performs semantic checks, validating register declarations, qubit indices, and classical variable bindings.

\textbf{Layer 3 - Intermediate Representation (IR):} Normalizes the AST into a hardware-independent \texttt{QuantumCircuit} instance. The IR maintains an ordered list of gate operations, classical bit registers, and condition dictionaries for feed-forward branching.

\textbf{Layer 4 - Decomposer Layer:} Lowers high-level or multi-controlled operations into elementary basis gates. For example, a 3-qubit Toffoli gate is decomposed into 6 CNOTs and 9 single-qubit rotations using Nielsen \& Chuang's canonical synthesis \cite{nielsen2010}.

\textbf{Layer 5 - Transpiler Layer:} Resolves hardware coupling constraints. It accepts a target topology graph (linear chain, 2D grid, or custom lattice) and uses Greedy BFS or SABRE \cite{li2019sabre} heuristics to insert SWAP gates, updating qubit mapping tables dynamically.

\textbf{Layer 6 - Simulator Layer:} Executes the physical circuit using one of two engines: exact statevector evolution using multi-dimensional NumPy array slicing (for $n \le 12$ qubits), or Matrix Product State tensor network contraction with SVD truncation (for $n \ge 20$ low-entanglement circuits) \cite{vidal2003mps}. An optional Kraus channel models physical noise.

\textbf{Layer 7 - Output Layer:} Collects simulation outcomes, formats statevector amplitudes, evaluates measurement probability distributions, generates Matplotlib and SVG circuit schematics, and syncs execution history to Supabase cloud storage.

\section{Activity Diagrams}
Activity diagrams capture the dynamic operational workflows of the system. We highlight the three central workflows in QVM: end-to-end circuit execution, transpilation routing, and simulator state evolution.

\subsection{Circuit Execution Flow}
Figure \ref{fig:activity_execution} depicts the complete lifecycle of a circuit execution, from loading source text through error validation, backend selection, simulation, and result presentation.

\begin{figure}[!htb]
\centering
\resizebox{0.88\textwidth}{!}{%
\begin{tikzpicture}[
    node distance=0.9cm and 1.2cm,
    startstop/.style={rectangle, rounded corners=8pt, minimum width=2.4cm, minimum height=0.65cm, text centered, draw=green!70!black, fill=green!15, font=\footnotesize\bfseries},
    process/.style={rectangle, rounded corners=3pt, minimum width=3.0cm, minimum height=0.7cm, text centered, draw=blue!75!black, fill=blue!8, font=\footnotesize, align=center},
    decision/.style={diamond, aspect=2, minimum width=2.2cm, minimum height=0.7cm, text centered, draw=orange!85!black, fill=orange!15, font=\tiny\bfseries, align=center},
    arrow/.style={-{Stealth[scale=0.95]}, line width=1pt, draw=black!80},
    errorstop/.style={rectangle, rounded corners=8pt, minimum width=2.4cm, minimum height=0.65cm, text centered, draw=red!70!black, fill=red!15, font=\footnotesize\bfseries}
]

\node[startstop] (start) {Start Execution};
\node[process, below=of start] (load) {Load Circuit File\\(QASM / JSON)};
\node[decision, below=of load] (checkfile) {File Exists \&\\Syntax Valid?};
\node[errorstop, right=2.0cm of checkfile] (err1) {Report Syntax Error\\Stop};

\node[process, below=of checkfile] (ir) {Build \texttt{QuantumCircuit}\\Intermediate Representation};
\node[decision, below=of ir] (transreq) {Transpilation\\Requested?};
\node[process, right=1.5cm of transreq] (transpile) {Decompose Gates \&\\Route on Target Graph};

\node[decision, below=of transreq] (qubitcheck) {Qubits $\le 12$?};
\node[process, below left=0.8cm and -0.5cm of qubitcheck] (statevec) {Select Statevector\\Simulator};
\node[decision, below right=0.8cm and -0.5cm of qubitcheck] (entangle) {Low Entanglement\\Circuit?};
\node[process, below=of entangle] (mps) {Select MPS\\Simulator (SVD)};
\node[errorstop, right=1.2cm of entangle] (err2) {Out of Memory\\Stop};

\node[process, below=2.6cm of qubitcheck] (init) {Initialize Quantum State\\$|0\dots 0\rangle$ and Memory};
\node[process, below=of init] (gates) {Evolve Quantum Gates\\\& Apply Born Collapse};
\node[process, below=of gates] (out) {Calculate Probabilities\\\& Render Histograms/Plots};
\node[startstop, below=of out] (end) {Return Results};

% Connections
\draw[arrow] (start) -- (load);
\draw[arrow] (load) -- (checkfile);
\draw[arrow] (checkfile) -- node[right, font=\tiny]{Yes} (ir);
\draw[arrow] (checkfile) -- node[above, font=\tiny]{No} (err1);

\draw[arrow] (ir) -- (transreq);
\draw[arrow] (transreq) -- node[above, font=\tiny]{Yes} (transpile);
\draw[arrow] (transreq) -- node[left, font=\tiny]{No} (qubitcheck);
\draw[arrow] (transpile) |- (qubitcheck);

\draw[arrow] (qubitcheck) -| node[above, font=\tiny]{Yes} (statevec);
\draw[arrow] (qubitcheck) -| node[above, font=\tiny]{No} (entangle);
\draw[arrow] (entangle) -- node[left, font=\tiny]{Yes} (mps);
\draw[arrow] (entangle) -- node[above, font=\tiny]{No} (err2);

\draw[arrow] (statevec) |- (init);
\draw[arrow] (mps) |- (init);
\draw[arrow] (init) -- (gates);
\draw[arrow] (gates) -- (out);
\draw[arrow] (out) -- (end);

\end{tikzpicture}%
}
\caption{Activity Diagram --- Circuit Execution Workflow}
\label{fig:activity_execution}
\end{figure}

\subsection{Transpilation Process}
Figure \ref{fig:activity_transpile} details the transpilation workflow. When a logical circuit is compiled to a physical chip topology, the transpiler decomposes high-level gates into basis primitives, queries the hardware coupling graph, and executes lookahead routing to insert SWAP gates.

\begin{figure}[!htb]
\centering
\resizebox{0.88\textwidth}{!}{%
\begin{tikzpicture}[
    node distance=0.85cm and 1.2cm,
    startstop/.style={rectangle, rounded corners=8pt, minimum width=2.4cm, minimum height=0.65cm, text centered, draw=green!70!black, fill=green!15, font=\footnotesize\bfseries},
    process/.style={rectangle, rounded corners=3pt, minimum width=3.2cm, minimum height=0.7cm, text centered, draw=blue!75!black, fill=blue!8, font=\footnotesize, align=center},
    decision/.style={diamond, aspect=2, minimum width=2.2cm, minimum height=0.7cm, text centered, draw=orange!85!black, fill=orange!15, font=\tiny\bfseries, align=center},
    arrow/.style={-{Stealth[scale=0.95]}, line width=1pt, draw=black!80}
]

\node[startstop] (start) {Start Transpilation};
\node[process, below=of start] (decomp) {Decompose Non-Native Gates\\(Toffoli, Controlled-$U$)};
\node[process, below=of decomp] (map) {Initialize Qubit Mapping\\$\text{logical} \leftrightarrow \text{physical}$};
\node[decision, below=of map] (strategy) {Select Routing Strategy};

\node[process, left=1.2cm of strategy] (greedy) {Greedy BFS Routing\\(Shortest Path)};
\node[process, right=1.2cm of strategy] (sabre) {SABRE Lookahead\\($W = \sum \text{decay}^i \cdot \text{dist}$)};

\node[decision, below=1.8cm of strategy] (adjacent) {Are Gate Qubits\\Physically Connected?};
\node[process, right=1.8cm of adjacent] (swap) {Evaluate Best SWAP\\Insert into Circuit \& Update Map};
\node[process, below=of adjacent] (emit) {Emit 2-Qubit Gate to\\Physical Circuit Stream};

\node[decision, below=of emit] (more) {More 2-Qubit Gates\\in Ready Queue?};
\node[decision, below=of more] (restore) {Restore Original\\Mapping?};
\node[process, right=1.8cm of restore] (revswap) {Insert Reverse SWAPs\\to Reset Qubit Index};
\node[startstop, below=of restore] (end) {Return Transpiled Circuit};

% Connections
\draw[arrow] (start) -- (decomp);
\draw[arrow] (decomp) -- (map);
\draw[arrow] (map) -- (strategy);

\draw[arrow] (strategy) -| node[above, font=\tiny]{Greedy} (greedy);
\draw[arrow] (strategy) -| node[above, font=\tiny]{SABRE} (sabre);
\draw[arrow] (greedy) |- (adjacent);
\draw[arrow] (sabre) |- (adjacent);

\draw[arrow] (adjacent) -- node[right, font=\tiny]{Yes} (emit);
\draw[arrow] (adjacent) -- node[above, font=\tiny]{No} (swap);
\draw[arrow] (swap) |- (adjacent);

\draw[arrow] (emit) -- (more);
\draw[arrow] (more) -- node[right, font=\tiny]{No} (restore);
\draw[arrow] (more.west) -- ++(-0.8,0) |- (adjacent.west) node[near start, left, font=\tiny]{Yes};

\draw[arrow] (restore) -- node[right, font=\tiny]{No} (end);
\draw[arrow] (restore) -- node[above, font=\tiny]{Yes} (revswap);
\draw[arrow] (revswap) |- (end);

\end{tikzpicture}%
}
\caption{Activity Diagram --- Transpilation and Routing Workflow}
\label{fig:activity_transpile}
\end{figure}

\subsection{Simulation Process}
Figure \ref{fig:activity_simulate} shows the internal flow of the simulation engine. The simulator initializes the quantum state $|\psi\rangle$ to $|0\dots 0\rangle$ and classical memory to zero. It iterates through the operation stream, applying unitary transformations, evaluating mid-circuit measurements via Born rule collapse, and processing classical conditionals.

\begin{figure}[!htb]
\centering
\resizebox{0.85\textwidth}{!}{%
\begin{tikzpicture}[
    node distance=0.85cm and 1.2cm,
    startstop/.style={rectangle, rounded corners=8pt, minimum width=2.4cm, minimum height=0.65cm, text centered, draw=green!70!black, fill=green!15, font=\footnotesize\bfseries},
    process/.style={rectangle, rounded corners=3pt, minimum width=3.2cm, minimum height=0.7cm, text centered, draw=blue!75!black, fill=blue!8, font=\footnotesize, align=center},
    decision/.style={diamond, aspect=2, minimum width=2.2cm, minimum height=0.7cm, text centered, draw=orange!85!black, fill=orange!15, font=\tiny\bfseries, align=center},
    arrow/.style={-{Stealth[scale=0.95]}, line width=1pt, draw=black!80}
]

\node[startstop] (start) {Start Simulator};
\node[process, below=of start] (init) {Initialize State $|0\dots 0\rangle$\\Reset Classical Memory, $PC \leftarrow 0$};
\node[decision, below=of init] (checkop) {Operations Remaining\\in Circuit?};

\node[decision, below=of checkop] (optype) {Operation Type?};

\node[process, left=1.2cm of optype] (gate) {Unitary Gate\\Tensor Slice Multiplication};
\node[process, right=1.2cm of optype] (meas) {Measurement\\Born Rule Sampling \& Collapse};
\node[process, below=1.2cm of optype] (ctrl) {Classical Condition\\Evaluate Bit Register};

\node[process, below=2.6cm of optype] (sync) {Advance Program Counter\\Update Classical Memory};
\node[startstop, right=2.5cm of checkop] (end) {Compute Probabilities\\Return Final State};

% Connections
\draw[arrow] (start) -- (init);
\draw[arrow] (init) -- (checkop);
\draw[arrow] (checkop) -- node[right, font=\tiny]{Yes} (optype);
\draw[arrow] (checkop) -- node[above, font=\tiny]{No} (end);

\draw[arrow] (optype) -| node[above, font=\tiny]{Gate} (gate);
\draw[arrow] (optype) -| node[above, font=\tiny]{Measure} (meas);
\draw[arrow] (optype) -- node[right, font=\tiny]{Control} (ctrl);

\draw[arrow] (gate) |- (sync);
\draw[arrow] (meas) |- (sync);
\draw[arrow] (ctrl) -- (sync);

\draw[arrow] (sync.west) -- ++(-1.0,0) |- (checkop.west);

\end{tikzpicture}%
}
\caption{Activity Diagram --- Simulator Execution Engine}
\label{fig:activity_simulate}
\end{figure}

\section{Class Diagram}
The static structure of QVM is captured in the class diagram shown in Figure \ref{fig:class}. The architecture reflects object-oriented encapsulation: the circuit is represented by \texttt{QuantumCircuit}, compilation by \texttt{Transpiler}, and simulation by \texttt{Simulator} and \texttt{MPSSimulator}.

\begin{figure}[!htb]
\centering
\resizebox{0.95\textwidth}{!}{%
\begin{tikzpicture}[
    cls/.style={
        rectangle split,
        rectangle split parts=3,
        draw=blue!85!black,
        rounded corners=2pt,
        line width=0.8pt,
        font=\scriptsize,
        align=left,
        inner sep=4pt,
        rectangle split part fill={blue!15, blue!4, white}
    },
    rel/.style={-{Stealth[scale=0.9]}, line width=0.9pt, draw=black!80},
    dep/.style={-{Stealth[scale=0.9]}, dashed, line width=0.9pt, draw=blue!75!black},
    comp/.style={{Diamond[open, length=6pt]}-{Stealth[scale=0.9]}, line width=0.9pt, draw=black!80}
]

% Center: QuantumCircuit (IR)
\node[cls, text width=5.2cm] (QC) at (0, 0) {
    \centering\textbf{QuantumCircuit}
    \nodepart{two}
    - num\_qubits: int\\
    - operations: List[Dict]\\
    - classical\_registers: Dict[str, int]
    \nodepart{three}
    + add\_operation(gate, qubits, params)\\
    + add\_classical\_register(name, size)\\
    + depth(): int \quad $\bullet$ \quad count(): int
};

% Top: CLI and API Interfaces
\node[cls, text width=4.6cm] (CLI) at (-3.2, 4.2) {
    \centering\textbf{CLI} \textit{(cli.py)}
    \nodepart{two}
    - args: argparse.Namespace
    \nodepart{three}
    + main(args: List[str])\\
    + execute\_circuit(path, opts)
};

\node[cls, text width=4.8cm] (API) at (3.2, 4.2) {
    \centering\textbf{APIServer} \textit{(api/app.py)}
    \nodepart{two}
    - app: FastAPI\\
    - supabase: Client | None
    \nodepart{three}
    + run(req: RunRequest): RunResponse\\
    + list\_circuits(): List[Dict]
};

% Left: Parsers
\node[cls, text width=4.6cm] (PARSER) at (-6.5, 0.8) {
    \centering\textbf{OpenQASM3Parser}
    \nodepart{two}
    - grammar: str (Lark LALR)\\
    - qubit\_map: Dict[str, int]
    \nodepart{three}
    + parse(text: str): QuantumCircuit\\
    - \_process\_ast(tree, cond)
};

\node[cls, text width=4.6cm] (DECOMP) at (-6.5, -3.2) {
    \centering\textbf{Decomposer}
    \nodepart{two}
    - native\_basis: List[str]
    \nodepart{three}
    + decompose(circuit): QuantumCircuit\\
    - \_synthesize\_toffoli(c1, c2, tgt)
};

% Mid-Left: Transpiler & TargetArchitecture
\node[cls, text width=5.2cm] (TRANS) at (-0.5, -4.5) {
    \centering\textbf{Transpiler}
    \nodepart{two}
    - architecture: TargetArchitecture\\
    - strategy: Literal['greedy', 'sabre']\\
    - restore\_mapping: bool
    \nodepart{three}
    + transpile(qc): QuantumCircuit\\
    - \_transpile\_sabre(qc): QuantumCircuit\\
    - \_heuristic\_cost(ready, map, decay)
};

\node[cls, text width=4.8cm] (ARCH) at (-0.5, -8.2) {
    \centering\textbf{TargetArchitecture}
    \nodepart{two}
    - num\_qubits: int\\
    - coupling\_graph: List[Tuple[int, int]]
    \nodepart{three}
    + is\_connected(q1, q2): bool\\
    + shortest\_path(u, v): List[int]
};

% Right: Simulator & MPSSimulator
\node[cls, text width=5.0cm] (SIM) at (6.5, 1.2) {
    \centering\textbf{Simulator} \textit{(Statevector)}
    \nodepart{two}
    - state: ndarray ($2^n$ complex128)\\
    - classical\_memory: Dict[str, List]\\
    - noise\_model: NoiseModel | None
    \nodepart{three}
    + simulate(circuit, seed): (ndarray, Dict)\\
    - \_apply\_single\_gate(state, U, k)\\
    - \_measure\_and\_collapse(qubit, rng)
};

\node[cls, text width=5.0cm] (MPS) at (6.5, -3.2) {
    \centering\textbf{MPSSimulator} \textit{(Tensor Train)}
    \nodepart{two}
    - chi\_max: int (Bond Limit)\\
    - tensors: List[ndarray]\\
    - svd\_threshold: float = $10^{-12}$
    \nodepart{three}
    + simulate(circuit): (ndarray, Dict)\\
    - \_contract\_two\_site(i, j)\\
    - \_truncate\_svd(theta, bond\_dim)
};

% Bottom Right: Visualizer
\node[cls, text width=4.8cm] (VIS) at (6.5, -7.5) {
    \centering\textbf{Visualizer} \textit{(visual.py)}
    \nodepart{two}
    - backend: str = 'Agg'
    \nodepart{three}
    + plot\_circuit(qc, output\_file)\\
    + plot\_histogram(counts, title)
};

% Connections / Relationships
\draw[dep] (CLI) -- node[right, font=\tiny]{invokes} (QC);
\draw[dep] (API) -- node[left, font=\tiny]{creates} (QC);
\draw[dep] (PARSER) -- node[above, font=\tiny]{generates} (QC);

\draw[rel] (TRANS) -- node[right, font=\tiny]{transforms} (QC);
\draw[comp] (TRANS) -- node[right, font=\tiny]{uses} (ARCH);
\draw[dep] (TRANS) -- (DECOMP);

\draw[rel] (SIM) -- node[above, font=\tiny]{executes} (QC);
\draw[rel] (MPS) -- node[above, font=\tiny]{executes} (QC);

\draw[dep] (CLI.west) -| (-6.5, 2.5) -- (PARSER.north);
\draw[dep] (CLI) |- (TRANS);
\draw[dep] (API) |- (TRANS);
\draw[dep] (CLI.east) -| (SIM.north);
\draw[dep] (API.east) -| (SIM.north);

\draw[dep] (SIM) -- (VIS);
\draw[dep] (MPS) -- (VIS);

\end{tikzpicture}%
}
\caption{QVM System Static Class Diagram}
\label{fig:class}
\end{figure}

\textbf{Core Class Descriptions:}
\begin{itemize}[leftmargin=*]
    \item \textbf{QuantumCircuit (\texttt{src/qvm/ir.py}):} The core data structure. Stores qubit counts, sequential gate operation dictionaries, and classical register sizes. Exposes helper methods: \texttt{add\_operation()}, \texttt{add\_classical\_register()}, and string serialization.
    \item \textbf{Simulator (\texttt{src/qvm/simulator.py}):} Manages dense statevector evolution. Represents states as NumPy complex arrays of shape $(2^n,)$ or $(2,2,\dots,2)$. Computes Born rule collapse upon measurement and maintains classical memory bit arrays.
    \item \textbf{MPSSimulator (\texttt{src/qvm/mps\_simulator.py}):} Represents states as a 1D tensor chain. Applies gates via two-site tensor contractions and truncates bond dimensions via SVD.
    \item \textbf{Transpiler (\texttt{src/qvm/transpiler.py}):} Directs physical qubit routing. Holds references to target coupling topologies and implements \texttt{\_transpile\_greedy()} and \texttt{\_transpile\_sabre()}.
    \item \textbf{TargetArchitecture (\texttt{src/qvm/architecture.py}):} Models physical coupling constraints as an undirected graph with adjacency checking and BFS shortest-path caching.
    \item \textbf{OpenQASM3Parser (\texttt{src/qvm/qasm3\_parser.py}):} Drives Lark LALR(1) parsing, building \texttt{QuantumCircuit} instances from OpenQASM 3.0 source code.
    \item \textbf{Decomposer (\texttt{src/qvm/decomposer.py}):} Expands high-level multi-qubit gates into native basis sets.
    \item \textbf{Visualizer (\texttt{src/qvm/visual.py}):} Renders Matplotlib circuit wire diagrams and probability histograms.
    \item \textbf{APIServer (\texttt{api/app.py}):} Exposes FastAPI REST endpoints (\texttt{POST /run}) with Pydantic payload validation and Supabase cloud persistence.
\end{itemize}

\section{Sequence Diagrams}
Sequence diagrams illustrate how objects collaborate over time during specific execution scenarios. We highlight three core interactions: command-line execution, web API request processing, and detailed transpilation routing.

\subsection{CLI Execution Flow}
Figure \ref{fig:seq_cli} models the interaction sequence when a user executes a quantum circuit from the terminal.

\begin{figure}[!htb]
\centering
\resizebox{0.92\textwidth}{!}{%
\begin{tikzpicture}[
  lifeline/.style={draw=black!50, dashed, line width=0.8pt},
  box/.style={rectangle, draw=blue!80!black, fill=blue!10, font=\footnotesize\bfseries, minimum height=0.7cm, inner sep=4pt},
  msg/.style={-{Stealth[scale=0.9]}, line width=0.9pt, draw=black!85},
  ret/.style={-{Stealth[scale=0.9]}, dashed, line width=0.9pt, draw=blue!75!black},
  note/.style={rectangle, draw=orange!70!black, fill=orange!10, font=\tiny, inner sep=3pt}
]

% Participants
\node[box] (user) at (0, 0) {User / Shell};
\node[box] (cli) at (2.5, 0) {CLI Handler};
\node[box] (parser) at (5.0, 0) {Parser};
\node[box] (trans) at (7.5, 0) {Transpiler};
\node[box] (sim) at (10.0, 0) {Simulator};
\node[box] (vis) at (12.2, 0) {Visualizer};

% Lifelines
\draw[lifeline] (user) -- (0, -6.0);
\draw[lifeline] (cli) -- (2.5, -6.0);
\draw[lifeline] (parser) -- (5.0, -6.0);
\draw[lifeline] (trans) -- (7.5, -6.0);
\draw[lifeline] (sim) -- (10.0, -6.0);
\draw[lifeline] (vis) -- (12.2, -6.0);

% Messages
\draw[msg] (0, -0.8) -- node[above, font=\tiny\bfseries]{invoke(file, --transpile, --sabre)} (2.5, -0.8);
\draw[msg] (2.5, -1.4) -- node[above, font=\tiny]{parse\_file(path)} (5.0, -1.4);
\draw[ret] (5.0, -2.0) -- node[above, font=\tiny]{QuantumCircuit IR} (2.5, -2.0);

\draw[msg] (2.5, -2.6) -- node[above, font=\tiny]{transpile(ir, target\_arch)} (7.5, -2.6);
\draw[ret] (7.5, -3.3) -- node[above, font=\tiny]{Physical Circuit + SWAPs} (2.5, -3.3);

\draw[msg] (2.5, -3.8) -- node[above, font=\tiny]{simulate(phys\_circuit)} (10.0, -3.8);
\draw[ret] (10.0, -4.5) -- node[above, font=\tiny]{Statevector, Memory, Counts} (2.5, -4.5);

\draw[msg] (2.5, -5.0) -- node[above, font=\tiny]{plot\_results(state)} (12.2, -5.0);
\draw[ret] (12.2, -5.4) -- node[above, font=\tiny]{Render Window / File} (2.5, -5.4);

\draw[ret] (2.5, -5.8) -- node[above, font=\tiny\bfseries]{Terminal Output Table} (0, -5.8);

\end{tikzpicture}%
}
\caption{Sequence Diagram --- CLI Execution Lifecycle}
\label{fig:seq_cli}
\end{figure}

\subsection{API Request Flow}
Figure \ref{fig:seq_api} illustrates the request-response cycle when an external client or web dashboard submits a quantum circuit through the FastAPI REST service.

\begin{figure}[!htb]
\centering
\resizebox{0.92\textwidth}{!}{%
\begin{tikzpicture}[
  lifeline/.style={draw=black!50, dashed, line width=0.8pt},
  box/.style={rectangle, draw=blue!80!black, fill=blue!10, font=\footnotesize\bfseries, minimum height=0.7cm, inner sep=4pt},
  msg/.style={-{Stealth[scale=0.9]}, line width=0.9pt, draw=black!85},
  ret/.style={-{Stealth[scale=0.9]}, dashed, line width=0.9pt, draw=blue!75!black}
]

% Participants
\node[box] (client) at (0, 0) {Web Client};
\node[box] (api) at (3.0, 0) {FastAPI Router};
\node[box] (pydantic) at (6.0, 0) {Pydantic Schema};
\node[box] (core) at (9.0, 0) {QVM Pipeline};
\node[box] (db) at (12.0, 0) {Supabase DB};

% Lifelines
\draw[lifeline] (client) -- (0, -5.2);
\draw[lifeline] (api) -- (3.0, -5.2);
\draw[lifeline] (pydantic) -- (6.0, -5.2);
\draw[lifeline] (core) -- (9.0, -5.2);
\draw[lifeline] (db) -- (12.0, -5.2);

% Messages
\draw[msg] (0, -0.8) -- node[above, font=\tiny\bfseries]{POST /run (JSON Payload)} (3.0, -0.8);
\draw[msg] (3.0, -1.4) -- node[above, font=\tiny]{validate\_request(body)} (6.0, -1.4);
\draw[ret] (6.0, -1.9) -- node[above, font=\tiny]{Validated Circuit Model} (3.0, -1.9);

\draw[msg] (3.0, -2.5) -- node[above, font=\tiny]{run\_simulation(circuit, shots, noise)} (9.0, -2.5);
\draw[ret] (9.0, -3.3) -- node[above, font=\tiny]{Amplitudes, Probabilities, Duration} (3.0, -3.3);

\draw[msg] (3.0, -3.8) -- node[above, font=\tiny]{save\_history(run\_record)} (12.0, -3.8);
\draw[ret] (12.0, -4.3) -- node[above, font=\tiny]{Record ID Stored} (3.0, -4.3);

\draw[ret] (3.0, -4.8) -- node[above, font=\tiny\bfseries]{HTTP 200 (JSON Results)} (0, -4.8);

\end{tikzpicture}%
}
\caption{Sequence Diagram --- Web API Request Lifecycle}
\label{fig:seq_api}
\end{figure}

\subsection{Transpilation Detail}
Figure \ref{fig:seq_transpile} details the interactions during circuit routing. The transpiler queries the hardware coupling graph, evaluates candidate SWAP operations, and updates the forward/inverse mapping tables.

\begin{figure}[!htb]
\centering
\resizebox{0.92\textwidth}{!}{%
\begin{tikzpicture}[
  lifeline/.style={draw=black!50, dashed, line width=0.8pt},
  box/.style={rectangle, draw=blue!80!black, fill=blue!10, font=\footnotesize\bfseries, minimum height=0.7cm, inner sep=4pt},
  msg/.style={-{Stealth[scale=0.9]}, line width=0.9pt, draw=black!85},
  ret/.style={-{Stealth[scale=0.9]}, dashed, line width=0.9pt, draw=blue!75!black}
]

% Participants
\node[box] (trans) at (0, 0) {Transpiler};
\node[box] (decomp) at (3.2, 0) {Decomposer};
\node[box] (arch) at (6.4, 0) {TargetArchitecture};
\node[box] (sabre) at (9.6, 0) {SABRE Evaluator};
\node[box] (outcir) at (12.5, 0) {Physical IR};

% Lifelines
\draw[lifeline] (trans) -- (0, -5.5);
\draw[lifeline] (decomp) -- (3.2, -5.5);
\draw[lifeline] (arch) -- (6.4, -5.5);
\draw[lifeline] (sabre) -- (9.6, -5.5);
\draw[lifeline] (outcir) -- (12.5, -5.5);

% Messages
\draw[msg] (0, -0.8) -- node[above, font=\tiny]{decompose\_circuit(ir)} (3.2, -0.8);
\draw[ret] (3.2, -1.4) -- node[above, font=\tiny]{Primitive Basis Circuit} (0, -1.4);

\draw[msg] (0, -2.0) -- node[above, font=\tiny]{is\_connected(p1, p2)} (6.4, -2.0);
\draw[ret] (6.4, -2.5) -- node[above, font=\tiny]{False (Non-Adjacent)} (0, -2.5);

\draw[msg] (0, -3.0) -- node[above, font=\tiny]{evaluate\_swaps(ready\_queue, map)} (9.6, -3.0);
\draw[ret] (9.6, -3.6) -- node[above, font=\tiny]{Best Candidate SWAP $(s_1, s_2)$} (0, -3.6);

\draw[msg] (0, -4.2) -- node[above, font=\tiny]{append\_swap(s1, s2)} (12.5, -4.2);
\draw[msg] (0, -4.7) -- node[above, font=\tiny]{append\_gate(p1, p2)} (12.5, -4.7);
\draw[ret] (12.5, -5.2) -- node[above, font=\tiny]{Physical Circuit Ready} (0, -5.2);

\end{tikzpicture}%
}
\caption{Sequence Diagram --- Transpiler Routing and SWAP Synthesis}
\label{fig:seq_transpile}
\end{figure}

\section{State Transition Diagrams}
Figure \ref{fig:state} shows the finite-state lifecycle of a quantum circuit as it passes through the QVM compilation and simulation pipeline.

\begin{figure}[!htb]
\centering
\resizebox{0.88\textwidth}{!}{%
\begin{tikzpicture}[
    node distance=1.3cm and 1.8cm,
    state/.style={rectangle, rounded corners=6pt, draw=blue!80!black, fill=blue!10, font=\footnotesize\bfseries, minimum height=0.8cm, minimum width=2.6cm, align=center},
    errstate/.style={rectangle, rounded corners=6pt, draw=red!80!black, fill=red!12, font=\footnotesize\bfseries, minimum height=0.8cm, minimum width=2.6cm, align=center},
    trans/.style={-{Stealth[scale=0.95]}, line width=1pt, draw=black!80}
]

\node[state] (raw) {Raw Text / File};
\node[state, below=of raw] (parsed) {Parsed IR};
\node[state, below=of parsed] (decomp) {Decomposed Circuit};
\node[state, below=of decomp] (transpiled) {Transpiled / Physical};
\node[state, below=of transpiled] (executing) {Executing State};
\node[state, below=of executing] (completed) {Completed State};

\node[errstate, right=2.5cm of parsed] (err) {Error State\\(Halt \& Log)};

% Transitions
\draw[trans] (raw) -- node[right, font=\tiny]{Lark Grammar Parse} (parsed);
\draw[trans] (parsed) -- node[right, font=\tiny]{Lower Gates to Basis} (decomp);
\draw[trans] (decomp) -- node[right, font=\tiny]{Route on Hardware Graph} (transpiled);
\draw[trans] (transpiled) -- node[right, font=\tiny]{Initialize Simulator} (executing);
\draw[trans] (executing) -- node[right, font=\tiny]{Evolve \& Measure} (completed);

% Error Transitions
\draw[trans] (raw.east) -| node[above, font=\tiny]{File Not Found} (err);
\draw[trans] (parsed.east) -- node[above, font=\tiny]{Syntax / Index Error} (err);
\draw[trans] (transpiled.east) -| node[above, font=\tiny]{Unroutable Graph} (err);
\draw[trans] (executing.east) -| node[above, font=\tiny]{Loop / Out-of-Memory} (err);

\end{tikzpicture}%
}
\caption{State Transition Diagram --- Quantum Circuit Lifecycle}
\label{fig:state}
\end{figure}

\section{Data Flow Diagrams}

\subsection{Level 1 DFD - System Level}
Figure \ref{fig:dfd1} illustrates the high-level data flows between external entities and the primary QVM execution process.

\begin{figure}[!htb]
\centering
\resizebox{0.88\textwidth}{!}{%
\begin{tikzpicture}[
    entity/.style={rectangle, draw=black!80, fill=gray!12, line width=1pt, minimum width=2.6cm, minimum height=0.9cm, font=\small\bfseries, align=center},
    process/.style={circle, draw=blue!80!black, fill=blue!10, line width=1.2pt, minimum size=2.8cm, font=\small\bfseries, align=center},
    datastore/.style={rectangle split, rectangle split parts=2, rectangle split horizontal, draw=black!80, fill=white, line width=0.8pt, minimum height=0.8cm, font=\footnotesize, align=center},
    flow/.style={-{Stealth[scale=0.95]}, line width=1pt, draw=black!75}
]

\node[entity] (user) at (-5.0, 1.5) {User / Client};
\node[process] (qvm) at (0, 0) {1.0\\QVM\\Runtime};
\node[entity] (fs) at (5.0, 1.5) {Local File\\System};
\node[entity] (cloud) at (5.0, -1.5) {Supabase\\Cloud DB};

\draw[flow] (user) -- node[above, font=\tiny\bfseries, sloped]{Circuit File / Config} (qvm);
\draw[flow] (qvm) -- node[below, font=\tiny\bfseries, sloped]{Results \& Histograms} (user);

\draw[flow] (fs) -- node[above, font=\tiny\bfseries, sloped]{Read .qasm / .json} (qvm);
\draw[flow] (qvm) -- node[below, font=\tiny\bfseries, sloped]{Write Plots / Exports} (fs);

\draw[flow] (qvm) -- node[above, font=\tiny\bfseries, sloped]{Save History} (cloud);
\draw[flow] (cloud) -- node[below, font=\tiny\bfseries, sloped]{Query Runs} (qvm);

\end{tikzpicture}%
}
\caption{Data Flow Diagram --- Level 1 System Level}
\label{fig:dfd1}
\end{figure}

\subsection{Level 2 DFD - Module Level}
Figure \ref{fig:dfd2} decomposes the central QVM runtime into its constituent data transformation processes.

\begin{figure}[!htb]
\centering
\resizebox{0.92\textwidth}{!}{%
\begin{tikzpicture}[
    proc/.style={rectangle, rounded corners=4pt, draw=blue!80!black, fill=blue!8, line width=1pt, minimum width=2.4cm, minimum height=0.8cm, font=\footnotesize\bfseries, align=center},
    store/.style={rectangle, draw=orange!80!black, fill=orange!8, line width=0.8pt, minimum width=2.2cm, minimum height=0.6cm, font=\tiny\bfseries, align=center},
    flow/.style={-{Stealth[scale=0.95]}, line width=1pt, draw=black!75}
]

\node[proc] (p1) at (-4.5, 2.0) {1.1 Parse\\Circuit Text};
\node[proc] (p2) at (-1.5, 2.0) {1.2 Decompose\\Gates};
\node[proc] (p3) at (1.5, 2.0) {1.3 Transpile\\\& Route};
\node[proc] (p4) at (4.5, 2.0) {1.4 Simulate\\Statevector/MPS};
\node[proc] (p5) at (1.5, -0.5) {1.5 Format\\Results};
\node[proc] (p6) at (-1.5, -0.5) {1.6 Export\\OpenQASM};

\node[store] (s1) at (1.5, 0.75) {Topology Graphs};
\node[store] (s2) at (-1.5, 0.75) {Basis Gate Sets};

\draw[flow] (p1) -- node[above, font=\tiny]{IR} (p2);
\draw[flow] (s2) -- (p2);
\draw[flow] (p2) -- node[above, font=\tiny]{Native IR} (p3);
\draw[flow] (s1) -- (p3);
\draw[flow] (p3) -- node[above, font=\tiny]{Physical IR} (p4);
\draw[flow] (p4) -- node[right, font=\tiny]{Amplitudes} (p5);
\draw[flow] (p3) |- (p6);

\end{tikzpicture}%
}
\caption{Data Flow Diagram --- Level 2 Module Level}
\label{fig:dfd2}
\end{figure}

\section{Entity-Relationship Diagram}
The Quantum Virtual Machine operates primarily as a computational runtime rather than a database-driven CRUD application. Core mathematical structures---statevectors, gate tensors, and classical bit arrays---live in volatile RAM during active simulation passes.

To support persistent circuit libraries, historical execution records, and cloud synchronization, QVM defines a relational data model shown in Figure \ref{fig:erd}. The persistence schema is deployed locally via SQLite/JSON and remotely via Supabase (PostgreSQL), tracking user accounts, parameterized circuits, hardware coupling topologies, execution runs, and numerical outcomes.

\begin{figure}[!htb]
\centering
\resizebox{0.95\textwidth}{!}{%
\begin{tikzpicture}[
    entity/.style={
        rectangle split,
        rectangle split parts=2,
        draw=blue!85!black,
        rounded corners=2pt,
        line width=0.8pt,
        font=\footnotesize,
        align=left,
        inner sep=4pt,
        rectangle split part fill={blue!15, white}
    },
    rel/.style={-{Stealth[scale=0.9]}, line width=0.9pt, draw=black!80},
    card/.style={font=\scriptsize\bfseries, text=blue!85!black}
]

% Entity 1: User
\node[entity, text width=4.8cm] (USER) at (-5.5, 3.2) {
    \centering\textbf{USER / DEVELOPER}
    \nodepart{two}
    \textbf{+ user\_id [PK]:} UUID\\
    - email: text\\
    - api\_key: varchar(64)\\
    - role: varchar(20)\\
    - created\_at: timestamptz
};

% Entity 2: Circuit
\node[entity, text width=5.2cm] (CIRCUIT) at (0, 3.2) {
    \centering\textbf{QUANTUM\_CIRCUIT}
    \nodepart{two}
    \textbf{+ circuit\_id [PK]:} UUID\\
    \textbf{\# user\_id [FK]:} UUID\\
    - name: varchar(100)\\
    - qasm\_source: text\\
    - num\_qubits: int\\
    - category: varchar(30)\\
    - created\_at: timestamptz
};

% Entity 3: Target Architecture
\node[entity, text width=5.0cm] (ARCH) at (5.8, 3.2) {
    \centering\textbf{TARGET\_ARCHITECTURE}
    \nodepart{two}
    \textbf{+ arch\_id [PK]:} varchar(32)\\
    - name: varchar(64)\\
    - num\_qubits: int\\
    - coupling\_map: jsonb\\
    - native\_basis: text[]\\
    - fidelity\_matrix: jsonb
};

% Entity 4: Execution Run
\node[entity, text width=5.5cm] (RUN) at (0, -1.8) {
    \centering\textbf{EXECUTION\_RUN}
    \nodepart{two}
    \textbf{+ run\_id [PK]:} UUID\\
    \textbf{\# circuit\_id [FK]:} UUID\\
    \textbf{\# arch\_id [FK]:} varchar(32)\\
    - engine: varchar(16) \textit{(sv / mps)}\\
    - routing: varchar(16) \textit{(greedy / sabre)}\\
    - shots: int \quad $\bullet$ \quad seed: bigint\\
    - noise\_config: jsonb\\
    - status: varchar(20)\\
    - submitted\_at: timestamptz
};

% Entity 5: Simulation Result
\node[entity, text width=5.4cm] (RESULT) at (5.8, -1.8) {
    \centering\textbf{SIMULATION\_RESULT}
    \nodepart{two}
    \textbf{+ result\_id [PK]:} UUID\\
    \textbf{\# run\_id [FK]:} UUID\\
    - probabilities: jsonb\\
    - sample\_counts: jsonb\\
    - expectation\_value: float\\
    - execution\_time\_ms: float\\
    - fidelity\_metric: float\\
    - completed\_at: timestamptz
};

% Relationships & Cardinalities
\draw[rel] (USER) -- node[above, card]{1} node[below, font=\tiny]{owns} (CIRCUIT) node[above, card, near end]{N};
\draw[rel] (CIRCUIT) -- node[right, card]{1} node[left, font=\tiny]{triggers} (RUN) node[right, card, near end]{N};
\draw[rel] (ARCH) |- node[above, card, near start]{1} node[below, font=\tiny, near start]{constrains} (RUN) node[right, card, near end]{N};
\draw[rel] (RUN) -- node[above, card]{1} node[below, font=\tiny]{yields} (RESULT) node[above, card, near end]{1};

\end{tikzpicture}%
}
\caption{Entity-Relationship Diagram --- Persistent Schema \& Execution Store}
\label{fig:erd}
\end{figure}

\begin{itemize}[leftmargin=*]
    \item \textbf{File-Based Storage:} Circuit definitions are serialized as standard \texttt{.qasm} or structured \texttt{.json} payloads. Simulation artifacts are stored as static PNG charts or raw numerical dumps.
    \item \textbf{Cloud Relational Schema:} Web sessions record executions using the five entities above, linking submitted QASM sources with compiled physical operations and output probability vectors for longitudinal auditing.
\end{itemize}

\section{Data Dictionary}
The following tables formally define the internal data structures used across the compilation and simulation pipeline.

\subsection{QuantumCircuit IR Structure}
Table \ref{tab:qc_struct} defines the core \texttt{QuantumCircuit} class structure.

\begin{table}[H]
\centering
\caption{QuantumCircuit Data Structure}
\label{tab:qc_struct}
\begin{tabular}{|p{3.5cm}|p{2.5cm}|p{7.5cm}|}
\hline
\textbf{Attribute} & \textbf{Data Type} & \textbf{Description} \\
\hline
\texttt{num\_qubits} & \texttt{int} & Total number of declared qubits in the circuit \\
\hline
\texttt{operations} & \texttt{list[dict]} & Ordered sequence of quantum and classical operations \\
\hline
\texttt{classical\_registers} & \texttt{dict[str, int]} & Map of classical register names to bit widths \\
\hline
\end{tabular}
\end{table}

\subsection{Operation Dictionary Structure}
Table \ref{tab:op_struct} specifies each gate and control dictionary stored inside \texttt{operations}.

\begin{table}[H]
\centering
\caption{Operation Data Structure}
\label{tab:op_struct}
\begin{tabular}{|p{3cm}|p{3cm}|p{7.5cm}|}
\hline
\textbf{Field} & \textbf{Data Type} & \textbf{Description} \\
\hline
\texttt{name} & \texttt{str} & Gate identifier (e.g., \texttt{h}, \texttt{cx}, \texttt{rx}, \texttt{measure}) \\
\hline
\texttt{qubits} & \texttt{list[int]} & Target physical/logical qubit indices (0-indexed) \\
\hline
\texttt{params} & \texttt{list[float]} & Rotation angles in radians (for parameterized gates) \\
\hline
\texttt{condition} & \texttt{dict | None} & Classical condition specification for branching \\
\hline
\texttt{target\_bit} & \texttt{tuple | None} & Destination register and index for measurement \\
\hline
\texttt{duration} & \texttt{str | None} & Timing directive (e.g., \texttt{"100ns"}) \\
\hline
\texttt{label} & \texttt{str | None} & Branch target label for control flow \\
\hline
\texttt{jump\_to} & \texttt{str | None} & Jump instruction target label \\
\hline
\texttt{classical\_op} & \texttt{dict | None} & Classical register assignment or boolean operation \\
\hline
\end{tabular}
\end{table}

\subsection{Condition Dictionary Structure}
Table \ref{tab:cond_struct} defines classical conditional branching criteria.

\begin{table}[H]
\centering
\caption{Condition Data Structure}
\label{tab:cond_struct}
\begin{tabular}{|p{3cm}|p{2.5cm}|p{8cm}|}
\hline
\textbf{Field} & \textbf{Data Type} & \textbf{Description} \\
\hline
\texttt{register} & \texttt{str} & Target classical register name \\
\hline
\texttt{index} & \texttt{int} & Specific bit index within register \\
\hline
\texttt{value} & \texttt{int} & Expected value ($0$ or $1$) required to fire operation \\
\hline
\end{tabular}
\end{table}

\subsection{Classical Operation Dictionary Structure}
Table \ref{tab:classic_struct} specifies classical bit operations.

\begin{table}[H]
\centering
\caption{Classical Operation Data Structure}
\label{tab:classic_struct}
\begin{tabular}{|p{3cm}|p{2.5cm}|p{8cm}|}
\hline
\textbf{Field} & \textbf{Data Type} & \textbf{Description} \\
\hline
\texttt{op} & \texttt{str} & Boolean operator (\texttt{=}, \texttt{\&}, \texttt{|}, \texttt{\^{}}, \texttt{\~{}}) \\
\hline
\texttt{target} & \texttt{tuple} & Destination tuple \texttt{(register\_name, bit\_index)} \\
\hline
\texttt{args} & \texttt{list} & Operand list (register references or integer literals) \\
\hline
\end{tabular}
\end{table}

\subsection{Target Architecture Structure}
Table \ref{tab:arch_struct} defines the hardware topology graph representation.

\begin{table}[H]
\centering
\caption{TargetArchitecture Data Structure}
\label{tab:arch_struct}
\begin{tabular}{|p{3cm}|p{2.5cm}|p{8cm}|}
\hline
\textbf{Attribute} & \textbf{Data Type} & \textbf{Description} \\
\hline
\texttt{num\_qubits} & \texttt{int} & Total number of physical qubits on the processor \\
\hline
\texttt{connectivity} & \texttt{list[tuple]} & Undirected edges \texttt{(q1, q2)} representing valid couplers \\
\hline
\texttt{name} & \texttt{str} & Architecture name (e.g., \texttt{"linear\_5"}, \texttt{"grid\_3x3"}) \\
\hline
\end{tabular}
\end{table}

\subsection{Simulation Results Structure}
Table \ref{tab:results_struct} defines the simulation output dictionary.

\begin{table}[H]
\centering
\caption{Simulation Results Data Structure}
\label{tab:results_struct}
\begin{tabular}{|p{3.5cm}|p{3.5cm}|p{6.5cm}|}
\hline
\textbf{Field} & \textbf{Data Type} & \textbf{Description} \\
\hline
\texttt{statevector} & \texttt{numpy.ndarray} & Complex amplitude array of dimension $2^n$ \\
\hline
\texttt{classical\_memory} & \texttt{dict[str, ndarray]} & Classical bit values stored by register \\
\hline
\texttt{probabilities} & \texttt{dict[str, float]} & Basis state bitstrings mapped to $|\alpha_i|^2$ \\
\hline
\texttt{execution\_time} & \texttt{float} & Simulation wall-clock duration in seconds \\
\hline
\end{tabular}
\end{table}

\section{Summary}
This chapter walked through the architectural blueprint of the Quantum Virtual Machine. We structured the compilation and runtime pipeline across seven distinct layers, applying the Strategy pattern so developers can swap routing heuristics without touching the execution core. Every legacy multi-page split diagram has been rebuilt directly as clean TikZ vector graphics---spanning the system context boundary, pipeline layers, execution flows, heuristic routing passes, simulator state transitions, static object models, and persistent relational schemas. These design decisions directly guided our implementation work in Chapter 5.
